\chapter{Repetition Test Library}
\label{ch:reptest}

This chapter discusses the library I created for benchmarking which was used in most experiments excepting any experiments utilizing FPGAs. Though not all
functionality was present through all experiments using it. It was inspired
and initially based architecturally on the one covered by Casey Muratori's Computer Enhance project \cite{muratori}. It builds off of and requires my own base
layer library of C code. It, like most of my C code, currently only supports Linux, but especially this particular library because it relies on the
Linux 'Perf' \cite{perf_2024} syscall API for collecting CPU hardware counters.

As one might assume from its title, it is primarily used to repeatedly test functions, collecting info
about them such as timing, page faults, in addition to the aforementioned hardware counters for things such as branch and cache misses. In my own biased opinion,
it provides this functionality in a straightforward and flexible way. It will continue to collect this info in a loop, collecting new minimums to try to
eliminate any OS-scheduling or other noise. It allows programs to own their own control flow, while simply marking up important
sections for benchmarking and data collection. It also allows users to provide additional info to calculate second order statistics such as throughput by
providing byte counts. Finally, it allows easy access to this data via options to output as CSV, containing both library generated data and user-provided data.
It also pushes live info as it is collected to the terminal.

\section{Motivation}
The experiments in this thesis required a library/harness with a particular combination of capabilities which none of the options I considered provided all of.

\paragraph{Hardware counters for diagnosis.}
The experiments in Chapter~\ref{ch:pact} compare operation counts against
runtime, and the gap between the two cannot be explained by operation
counts alone nor do they necessarily even predict real performance. At this point the library did not collect hardware counter events, but explaining
performance in reality requires cache misses, branch mispredictions,
and other hardware events. Collecting these through source-level
instrumentation, as I did with the counting macros in that chapter,
perturbs the performance being measured, in addition to being tedious to insert. Hardware counters read via
the Linux \texttt{perf} interface \cite{perf_2024} avoid this, so a harness that exposed them directly was necessary.

\paragraph{Measurement of marked regions, not whole processes.}
Tools such as \texttt{perf stat} report counters for entire processes, which includes setup, input generation, verification, etc. I wanted
measurements scoped to individual kernels, so the harness had to let the program and user mark the regions of interest.

\paragraph{Program-owned control flow.}
Many benchmarking frameworks invert control. Often they require the user to register a
function and the framework or library decides when and how to call it. This can be
awkward for kernels that need per-experiment setup, such as matrix
generation at a chosen density, or that we wish to be driven from a variable-sweep loop in
the program, or that have special blocking transformations that need to be done on input matrices.
I wanted the program to own its own control flow and simply mark the section(s) to be repeated and measured.

\paragraph{Simplicity.}
All the code for this thesis is written in C on top of a small base layer library of my own.
I considered heavier weight libraries like PAPI \cite{papi}, which would have potentially
allowed me to collect even more cross-microarchitecture (TODO: better way to say 'cross-microarchitecture'?) hardware counter events such as
flop/load/store instructions retired easier. However, I decided against it in the interest of keeping
benchmarking complexity low. I wanted to save most of my mental complexity budget for dealing with the actual experiments. As well, its always
more interesting to build your own thing than use a library.

\section{Specific Measurements}
This section defines the quantities reported in Chapters~\ref{ch:pact} and~\ref{ch:kron_taco}.

\paragraph{Time.}
Runtime is measured with the x86 timestamp counter (\texttt{RDTSC} instruction) and
reported in timestamp ``cycles.'' This counter advances at a fixed rate regardless of actual core frequency
scaling, so these units are ticks of a constant-rate clock, not cycles
of the core clock. Each kernel is repeated and the minimum observed time is kept, again to minimize scheduling noise, first touch page faults,
and cold-cache effects. The minimum is therefore a best-case, warm-state estimate.

\paragraph{Hardware counters.}
Hardware-counted events are read through the Linux \texttt{perf} interface
\cite{perf}, scoped to the timed region. Two generalized events are used.
\emph{Cache misses} corresponds to \texttt{PERF\_COUNT\_HW\_CACHE\_MISSES}, which
usually counts last-level cache misses, i.e., accesses that must go to main memory. \emph{Branch mispredictions} corresponds to
\texttt{PERF\_COUNT\_HW\_BRANCH\_MISSES}, the number of mispredicted branches.

\paragraph{Page faults.}
Memory returned by an allocator is often not immediately backed by physical pages.
The first access to each page raises a \emph{page fault}, which the operating
system remedies by mapping an actual page. Page faults therefore measure
first-touch cost: a kernel that writes freshly allocated memory incurs roughly one
fault per page touched (typically 4~KiB, 2~MiB, or 1~GiB), and each fault costs kernel time that is
included in the measured runtime. This matters for sparse formats because
their working-set sizes can scale with both matrix dimension and the number of non-zeros. As well, TACO uses the C standard library allocator malloc and calloc
which are not known to have the best behaviour regarding pre-faulting.

\paragraph{Allocations.}
In addition to counters, each run records the total bytes requested from the
allocators malloc and calloc. This is \emph{requested}, not actually mapped or touched memory,
so it can exceed what page faults show if memory is allocated but never written, as is often the case with TACO generated code.

\section{Example Usage}
Below I provide a simple example program demonstrating the library's use. Note that the user is free do define the order and sections of interest.
The example usage below is not the only way to use the library, but it has the general structure most of my experiments ended up in.

\note{Probably this giant ass listing needs to go in the appendix...}
\begin{lstlisting}[language=C, caption={Example usage of Repetition Test Library}]
// NOTE: Define the interface that all the functions we want to test conform to.
typedef void (*Test_Function)(u8 *buffer, usize buffer_size);
static
void read_buffer_linearly(u8 *buffer, usize buffer_size)
{
  ... // Omitted for brevity
}
static
void read_buffer_randomly(u8 *buffer, usize buffer_size)
{
  ... // Omitted for brevity
}

// NOTE: Each entry is just going to be the function and
// then a little string to name it.
typedef struct Function_Entry Function_Entry;
struct Function_Entry
{
  const char    *name;
  Test_Function function;
};
// NOTE: Ok now we can just make an array of all the functions we want to test.
Function_Entry entries[] =
{
  {"linearly", read_buffer_linearly},
  {"randomly", read_buffer_randomly},
};

int main(int argc, char **argv)
{
  Arena arena = arena_make(.reserve_size = GB(64));
  u64 cpu_timer_frequency = estimate_cpu_timer_freq();

  // NOTE: How many seconds to keep trying for a new min, will reset once
  // we find a new minimum.
  u32 seconds_to_try_for_min = 1;

  // NOTE: User may provide custom fields.
  Repetition_Series *series = repetition_series_make(60 * STATIC_COUNT(entries), ((const char *[]){"buffer_size", "function"}));

  for (usize buffer_size = MB(1); buffer_size <= MB(4); buffer_size *= 4)
  {
    for (usize test_func_idx = 0; test_func_idx < STATIC_COUNT(entries); test_func_idx++)
    {
      // NOTE: Allocating and deallocating every time to make it fair page-fault wise.
      u8 *buffer = os_allocate(buffer_size, OS_ALLOCATION_COMMIT);

      Function_Entry *entry = &entries[test_func_idx];

      Repetition_Tester tester = repetition_series_new_tester(series, buffer_size,
                                                              cpu_timer_frequency,
                                                              seconds_to_try_for_min,
                                                              "--- %s @ %lu MB ---", entry->name,
                                                              buffer_size / MB(1));

      // NOTE: User may set fields of current test.
      repetition_series_set_field(series, "buffer_size", "%lu", buffer_size);
      repetition_series_set_field(series, "function", entry->name);

      // NOTE: The main loop, it will continue testing until the minimum time has not decreased in user-provided time-frame.
      while (repetition_series_is_testing(series, &tester))
      {
        repetition_tester_begin_time(&tester);
        entry->function(buffer, buffer_size);
        repetition_tester_close_time(&tester);

        repetition_tester_count_bytes(&tester, buffer_size);
      }

      os_deallocate(buffer, buffer_size);
    }
  }

  // NOTE: Self-explanatory
  repetition_series_save_csv(series, "out.csv");
}
\end{lstlisting}
